% Shared setup for the four-condition, full-12-epoch Orin experiment.
% Evidence: ../orin_scheduling/results/feature_20261113_prefetch_full_v2/
\begingroup
\color{blue}
\subsection{Coexisting with Periodic Inference on Jetson Orin}
\label{subsec:orin-periodic}

We test whether reordering training work and computing features ahead of time can speed up background training while foreground inference runs regularly.
We compare four methods: FIFO, Reorder, FIFO with prefetching, and Reorder with prefetching.
FIFO without feature caching is the baseline.
Both optimizations run on the device and leave offline proxy generation and policy training unchanged.

\noindent \textbf{Setup.}
We use a Jetson AGX Orin developer kit with 64\,GB memory in its existing 30\,W power mode.
It runs PyTorch 2.0.0+nv23.05 and CUDA 11.4.
We use the released ResNet50 proxy networks, policy model, and classifier on Caltech-101.
The proxy networks stay frozen, and we train the classifier for 12 epochs with batch size four and an SGD learning rate of 0.001.
All methods use the same models and store tensors in FP32.
We turn off the Python TF32 settings, although the released TorchScript convolution operators still allow TF32 internally.

For the four-method comparison, a foreground inference request arrives every 105\,ms during training.
We measure its response time from its scheduled arrival until the GPU finishes it, including any waiting time.
We set a response-time target of 33.20\,ms before the experiments and process every request.

\noindent \textbf{Scheduling.}
We use the FIFO and Reorder rules in Section~\ref{subsec:reorder}.
All four methods give foreground requests priority and leave a 2\,ms safety margin before the next request.
Reorder looks ahead by up to 128 batches.
Without caching, it stops letting later batches pass the oldest batch after eight such passes.
With caching, only feature computation is reordered; the remaining training work keeps the original batch order.

We estimate how long each operation will take from measured runtimes.
The initial estimate is 1.15 times the 95th-percentile runtime measured without foreground inference.
We then update it to 1.10 times the 95th percentile of the latest 32 measurements, but never below the initial estimate.
For any operation that may update the classifier, we reserve enough time for a full backward update before deciding whether the loss calls for one.

\noindent \textbf{Timing.}
Each reported result comes from one run.
We use a fixed seed to shuffle the run order within each experiment.
Each run starts from the same classifier and routing state and uses the same random seed.
Training time covers all 12 training epochs, including waiting time and scheduling overhead, but excludes evaluation.
Total time also includes the 12 full evaluations and the time needed to finish any remaining foreground requests.
We load the models and prepared images into GPU memory before timing starts.
Initial policy inference, runtime measurements for the starting estimates, and warmup also take place before timing starts.
We check the logs to confirm that all training batches, evaluations, and foreground requests finish, that batches follow the required order, and that the recorded times are consistent.
\par
\endgroup
